\begin{sublemma}
\uses{score-def,factorization-def,balance-def}
\label{score-lower-3}
\lean{Erdos392.Factorization.lower_score_3}
% at /home/runner/work/trough/trough/library/PrimeNumberTheoremAnd/IEANTN/Erdos392.lean:834.0-864.51
\leanok
If there is a prime $p$ in deficit less than $L$, one can remove it without increasing the score.
\end{sublemma}
\begin{proof}
\uses{waste-def,score-eq,score-lower-2,score-lower-1}
\leanok
Without loss of generality we may assume that one is not in the previous two
  situations, i.e., wlog there are no surplus primes and all primes in deficit are at most $L$.
  If all deficit primes multiply to $n$ or less, add that product to the factorization (this
  increases the waste by at most $\log n$, but we save a $\log n$ from now being in balance).
  Otherwise, greedily multiply all primes together while staying below $n$ until one cannot do so
  any further; add this product to the factorization, increasing the waste by at most $\log L$.
\end{proof}
